\documentclass[10.5pt]{article}
\usepackage[letterpaper, margin=0.6in, top=0.5in, bottom=0.5in]{geometry}
\usepackage{amsmath,amssymb,amsthm}
\usepackage{fancyhdr}
\usepackage{xcolor}
\usepackage{tikz}
\usepackage{enumitem}
\usepackage{tcolorbox}
\usepackage{hyperref}

% Define colors matching instructor board annotations
\definecolor{chalkblue}{RGB}{0, 105, 185}
\definecolor{chalkred}{RGB}{205, 35, 35}
\definecolor{chalkpink}{RGB}{230, 90, 120}
\definecolor{chalkcyan}{RGB}{15, 160, 190}
\definecolor{chalkgreen}{RGB}{20, 130, 50}
\definecolor{boxbg}{RGB}{252, 252, 254}
\definecolor{boxframe}{RGB}{150, 160, 180}

\pagestyle{fancy}
\fancyhf{}
\renewcommand{\headrulewidth}{0pt}
\renewcommand{\footrulewidth}{0pt}

\setlength{\parindent}{0pt}
\setlength{\parskip}{2.5pt}

\begin{document}

% ==============================================================================
% PAGE 1
% ==============================================================================
\noindent \textbf{\large Math 105L} \hfill \textbf{\large Lecture 2-3}\\[1pt]
{\textbf{\Large Function Transformations: Shifts, Scalings, \& Roots}}\\[-5pt]
\rule{\textwidth}{1.2pt}

\vspace{0.03in}

\textbf{\large 1. Summary of Elementary Transformations}

Let $y = f(x)$ be a given parent function and $c > 0$ a positive constant:

\begin{tcolorbox}[colback=boxbg, colframe=boxframe, boxrule=0.8pt, arc=3pt, left=8pt, right=8pt, top=4pt, bottom=4pt]
\begin{itemize}[leftmargin=1.2em, topsep=1pt, itemsep=3pt]
  \item \textbf{Vertical Translations:}
  \begin{itemize}[label=--]
    \item $y = f(x) + c$: Shifts graph \textbf{up} by $c$ units.
    \item $y = f(x) - c$: Shifts graph \textbf{down} by $c$ units.
  \end{itemize}
  \item \textbf{Horizontal Translations:}
  \begin{itemize}[label=--]
    \item $y = f(x - c)$: Shifts graph \textbf{right} by $c$ units (replaces $x$ with $x-c$).
    \item $y = f(x + c)$: Shifts graph \textbf{left} by $c$ units (replaces $x$ with $x+c$).
  \end{itemize}
  \item \textbf{Scalings (Stretches and Compressions):}
  \begin{itemize}[label=--]
    \item $y = c \cdot f(x)$ ($c > 1$): Vertical stretch by factor $c$.
    \item $y = f(c x)$ ($c > 1$): \textbf{Horizontal compression} by factor $c$ (inputs are divided by $c$).
  \end{itemize}
\end{itemize}
\end{tcolorbox}

\vspace{0.03in}

\textbf{\large 2. Classroom Polynomial Case Study: Tracking Roots}

Consider the cubic polynomial with given roots at $x = -1, 1, 3$:
\[
  f(x) = (x + 1)(x - 1)(x - 3)
\]

\textbf{Task:} Construct the transformed function $g(x) = f\big(2(x - 1)\big)$, simplify its factors, and determine its new roots.

\begin{itemize}[leftmargin=1.2em, topsep=1pt, itemsep=4pt]
  \item \textbf{Step 1: Substitute the transformed input $2(x - 1)$ into each factor:}
  \begin{align*}
    g(x) &= \Big(\big[2(x - 1)\big] + 1\Big)\Big(\big[2(x - 1)\big] - 1\Big)\Big(\big[2(x - 1)\big] - 3\Big) \\
         &= \big(2x - 2 + 1\big)\big(2x - 2 - 1\big)\big(2x - 2 - 3\big) \\
         &= \mathbf{(2x - 1)(2x - 3)(2x - 5)}
  \end{align*}

  \item \textbf{Step 2: Solve for the new roots ($g(x) = 0$):}
  \begin{itemize}[label=--]
    \item $2x - 1 = 0 \implies \mathbf{x = \frac{1}{2}}$
    \item $2x - 3 = 0 \implies \mathbf{x = \frac{3}{2}}$
    \item $2x - 5 = 0 \implies \mathbf{x = \frac{5}{2}}$
  \end{itemize}
\end{itemize}

\vspace{0.04in}

\begin{tcolorbox}[colback=chalkblue!6, colframe=chalkblue, boxrule=1pt, arc=3pt, left=8pt, right=8pt, top=4pt, bottom=4pt]
\textbf{Geometric Transformation Verification:}\\[2pt]
The transformation $f\big(2(x - 1)\big)$ applies two horizontal operations to the original $x$-coordinates:
\begin{enumerate}[leftmargin=1.4em, topsep=1pt, itemsep=2pt]
  \item \textbf{Horizontal Compression by 2:} Divides each original root by $2$:
  \[
    -1 \to -0.5, \quad 1 \to 0.5, \quad 3 \to 1.5
  \]
  \item \textbf{Horizontal Shift Right by 1:} Adds $+1$ to each coordinate:
  \[
    -0.5 + 1 = \mathbf{0.5}, \quad 0.5 + 1 = \mathbf{1.5}, \quad 1.5 + 1 = \mathbf{2.5} \qquad \checkmark
  \]
\end{enumerate}
Algebra and geometry agree completely!
\end{tcolorbox}

\end{document}

